% ═══════════════════════════════════════════════════════════════
% SL-08d-ML: Stundenleistung MUSTERLÖSUNG
% ═══════════════════════════════════════════════════════════════

\documentclass[11pt, a4paper]{article}

\newif\ifswmodus
\swmodusfalse

\usepackage[utf8]{inputenc}
\usepackage[T1]{fontenc}
\usepackage[ngerman]{babel}
\usepackage{helvet}
\renewcommand{\familydefault}{\sfdefault}

\usepackage[margin=2cm]{geometry}
\usepackage{enumitem}
\usepackage{tabularx}
\usepackage{booktabs}

\usepackage{xcolor}
\usepackage{tcolorbox}
\tcbuselibrary{skins}

\usepackage{fancyhdr}
\usepackage{lastpage}
\usepackage{needspace}

\definecolor{spanisch-color}{HTML}{c41e3a}
\definecolor{grau-color}{HTML}{888888}
\definecolor{hellgrau-color}{HTML}{f5f5f5}
\definecolor{loesung-color}{HTML}{c62828}

\definecolor{sw-dark}{HTML}{000000}
\definecolor{sw-gray}{HTML}{666666}
\definecolor{sw-white}{HTML}{ffffff}

\ifswmodus
    \colorlet{spanisch}{sw-dark}
    \colorlet{grau}{sw-gray}
    \colorlet{hellgrau}{sw-white}
    \colorlet{loesung}{sw-dark}
\else
    \colorlet{spanisch}{spanisch-color}
    \colorlet{grau}{grau-color}
    \colorlet{hellgrau}{hellgrau-color}
    \colorlet{loesung}{loesung-color}
\fi

\newcommand{\fachfarbe}{spanisch}

\newcommand{\thema}{Stundenleistung: Vacaciones}
\newcommand{\fach}{Spanisch}
\newcommand{\klasse}{7}
\newcommand{\maxpunkte}{20}
\newcommand{\zeit}{25 Minuten}
\newcommand{\datum}{\today}

\pagestyle{fancy}
\fancyhf{}
\renewcommand{\headrulewidth}{0.4pt}
\renewcommand{\footrulewidth}{0.4pt}

\lhead{\fach{} -- Klasse \klasse}
\chead{\textcolor{loesung}{\textbf{SL-08d MUSTERLÖSUNG}}}
\rhead{\datum}

\lfoot{\zeit}
\cfoot{}
\rfoot{Seite \thepage\ von \pageref{LastPage}}

\newcommand{\punktebox}[1]{\hfill\fbox{\small \_/#1}}
\newcommand{\luecke}[1]{\rule{#1}{0.4pt}}
\newcommand{\lsg}[1]{\textcolor{loesung}{\textbf{#1}}}

\newtcolorbox{infobox}{
    colback=hellgrau,
    colframe=\fachfarbe,
    boxrule=1pt,
    arc=0pt,
    left=8pt, right=8pt, top=8pt, bottom=8pt
}

\newenvironment{aufgabe}[1]{%
    \needspace{5\baselineskip}%
    \nopagebreak[4]%
    \vspace{10pt}
    \noindent\textbf{\textcolor{\fachfarbe}{Aufgabe #1}}
    \nopagebreak[4]%
    \vspace{5pt}
    \nopagebreak[4]%
    \begin{list}{}{\leftmargin=0pt}
    \item[]
}{%
    \end{list}
}

\newcommand{\es}[1]{\textcolor{\fachfarbe}{\textbf{#1}}}

\begin{document}

\begin{center}
    {\Large\bfseries\textcolor{\fachfarbe}{\thema}}\\[0.3em]
    {\large\textcolor{loesung}{\textbf{--- MUSTERLÖSUNG ---}}}
\end{center}

\vspace{5pt}

\noindent
\begin{tabularx}{\textwidth}{|X|X|X|}
\hline
\textbf{Name:} \lsg{Musterschüler/in} &
\textbf{Klasse:} \klasse &
\textbf{Datum:} \datum \\
\hline
\end{tabularx}

\vspace{10pt}

\begin{infobox}
\textbf{Zeit:} \zeit{} | \textbf{Punkte:} \maxpunkte{} | \textbf{Hilfsmittel:} keine
\end{infobox}

\vspace{15pt}

% ═══════════════════════════════════════════════════════════════
% AUFGABE 1
% ═══════════════════════════════════════════════════════════════

\begin{aufgabe}{1 -- Reisevokabular zuordnen}
Verbinde das spanische Wort mit der deutschen Bedeutung. \punktebox{4}
\end{aufgabe}

\begin{center}
\begin{tabular}{rcl}
\es{el tren} & \lsg{---} & das Flugzeug \lsg{$\rightarrow$ el tren = der Zug} \\[0.8em]
\es{el avión} & \lsg{---} & der Strand \lsg{$\rightarrow$ el avión = das Flugzeug} \\[0.8em]
\es{la playa} & \lsg{---} & der Zug \lsg{$\rightarrow$ la playa = der Strand} \\[0.8em]
\es{el hotel} & \lsg{---} & das Hotel \lsg{$\rightarrow$ el hotel = das Hotel} \\
\end{tabular}
\end{center}

\textit{\textcolor{grau}{Bewertung: 1 Punkt pro richtige Zuordnung}}

\vspace{15pt}

% ═══════════════════════════════════════════════════════════════
% AUFGABE 2
% ═══════════════════════════════════════════════════════════════

\begin{aufgabe}{2 -- Wettersätze bilden}
Schreibe den passenden spanischen Satz. \punktebox{4}
\end{aufgabe}

\noindent
a) Es ist warm. $\rightarrow$ \lsg{Hace calor.}

\vspace{0.8em}

\noindent
b) Es regnet. $\rightarrow$ \lsg{Llueve.}

\vspace{0.8em}

\noindent
c) Es ist kalt. $\rightarrow$ \lsg{Hace frío.}

\vspace{0.8em}

\noindent
d) Die Sonne scheint. $\rightarrow$ \lsg{Hace sol.}

\textit{\textcolor{grau}{Bewertung: 1 Punkt pro richtige Antwort. Bei "Hace" + Nomen: 0,5P für richtige Struktur, 0,5P für richtiges Wort}}

\vspace{15pt}

% ═══════════════════════════════════════════════════════════════
% AUFGABE 3
% ═══════════════════════════════════════════════════════════════

\begin{aufgabe}{3 -- Pretérito: Verben konjugieren}
Setze die richtige Form des Pretérito indefinido ein. \punktebox{6}
\end{aufgabe}

\noindent
a) Yo \lsg{viajé} (viajar) a Italia el año pasado. \hfill (ich reiste)

\vspace{0.8em}

\noindent
b) Tú \lsg{visitaste} (visitar) Barcelona ayer. \hfill (du besuchtest)

\vspace{0.8em}

\noindent
c) Mi hermano \lsg{nadó} (nadar) en el mar. \hfill (er schwamm)

\vspace{0.8em}

\noindent
d) Nosotros \lsg{bailamos} (bailar) en la fiesta. \hfill (wir tanzten)

\vspace{0.8em}

\noindent
e) Vosotros \lsg{comprasteis} (comprar) recuerdos. \hfill (ihr kauftet)

\vspace{0.8em}

\noindent
f) Ellos \lsg{tomaron} (tomar) el sol. \hfill (sie sonnten sich)

\textit{\textcolor{grau}{Bewertung: 1 Punkt pro richtige Form. Fehlender Akzent bei viajé/nadó: nur 0,5P}}

\vspace{15pt}

% ═══════════════════════════════════════════════════════════════
% AUFGABE 4
% ═══════════════════════════════════════════════════════════════

\begin{aufgabe}{4 -- Urlaubsbericht schreiben}
Schreibe einen kurzen Bericht über einen Urlaub (4 Sätze im Pretérito). \punktebox{6}
\end{aufgabe}

\textbf{Beispiellösung:}

\lsg{El año pasado viajé a España con mi familia.}

\lsg{Visitamos Barcelona y la playa.}

\lsg{Nadé en el mar y tomé el sol.}

\lsg{Fue muy bonito y divertido.}

\vspace{0.5em}

\textit{\textcolor{grau}{Bewertungskriterien:}}
\begin{itemize}[nosep]
    \item \textit{\textcolor{grau}{4 Sätze im Pretérito: 4P (1P pro Satz mit korrekter Form)}}
    \item \textit{\textcolor{grau}{Passender Wortschatz (Reisen/Wetter): 1P}}
    \item \textit{\textcolor{grau}{Sprachliche Korrektheit: 1P}}
\end{itemize}

% ═══════════════════════════════════════════════════════════════
% PUNKTEÜBERSICHT
% ═══════════════════════════════════════════════════════════════

\vspace{25pt}
\hrule
\vspace{10pt}

\noindent
\begin{tabularx}{\textwidth}{|X|c|c|c|c||c|}
\hline
\textbf{Aufgabe} & 1 & 2 & 3 & 4 & \textbf{Gesamt} \\
\hline
\textbf{max. Punkte} & 4 & 4 & 6 & 6 & \textbf{20} \\
\hline
\textbf{erreicht} & \lsg{4} & \lsg{4} & \lsg{6} & \lsg{6} & \lsg{20} \\
\hline
\end{tabularx}

\vspace{10pt}

\begin{flushright}
\textbf{Note:} \lsg{1}
\end{flushright}

\vspace{1em}

\textbf{Notenschlüssel (Sek I, 14 NP):}

\begin{tabular}{|c|c|c|c|c|c|c|c|c|c|c|c|c|c|c|}
\hline
\textbf{NP} & 14 & 13 & 12 & 11 & 10 & 9 & 8 & 7 & 6 & 5 & 4 & 3 & 2 & 1 \\
\hline
\textbf{ab P.} & 20 & 19 & 18 & 17 & 16 & 15 & 13 & 12 & 11 & 9 & 8 & 7 & 5 & 4 \\
\hline
\end{tabular}

\vspace{1em}

\textbf{Alternativ Notenstufen:}

\begin{tabular}{|c|c|c|c|c|c|c|}
\hline
\textbf{Note} & 1 & 2 & 3 & 4 & 5 & 6 \\
\hline
\textbf{ab Punkte} & 19 & 16 & 12 & 8 & 4 & $<$4 \\
\hline
\end{tabular}

\end{document}
